% exp-fips203-mlkem-pke-stage3-routea-mod-vec Stage3 向量预研实现方案
% 编译：bash ../../../../scripts/xelatex-clean.sh exp-fips203-mlkem-pke-stage3-routea-mod-vec-实现方案-customspec.tex
\documentclass[UTF8,a4paper,11pt]{ctexart}
\usepackage{amsmath,amssymb}
\usepackage{booktabs}
\usepackage{geometry}
\usepackage{hyperref}
\usepackage{longtable}
\usepackage{array}
\usepackage{seqsplit}
\geometry{margin=2.2cm}
% 表格列内自动换行；长 API 名在 \texttt 中可断行（与 Stage1 customspec 一致）
\newcolumntype{L}[1]{>{\raggedright\arraybackslash}p{#1}}
\newcommand{\apiname}[1]{\texttt{\seqsplit{#1}}}
\hypersetup{colorlinks=true,linkcolor=blue,urlcolor=blue}

\title{exp-fips203-mlkem-pke-stage3-routea-mod-vec\\F203 Stage3 RouteA + mod $q$ 向量预研方案}
\author{ascendc 工程（预研）}
\date{2026-05-19}

\begin{document}
\maketitle

\begin{abstract}
本文档为 \texttt{exp-fips203-mlkem-pke-stage3-routea-mod-vec/} 的\textbf{计划书闭环}：
\texttt{mat\_c [16,512] int32}（Stage2 Cube 输出）$\rightarrow$ RouteA $\rightarrow$ Stage3.1 $\bmod\,q$ $\rightarrow$ \texttt{out [8,256] int32}。
在查阅 CANN 9.0.0 Ascend C API 与 dav\_c220 实现后，\textbf{Atlas A2 上无法对 RouteA 合并与 Stage3.1 做到 100\% 单向量而保持 int64/除法语义}；
本用例采用\textbf{向量 GM 搬运 + 标量 int64 合并与 mod} 的混合实现，并在下表明确标量缺口与备选 API 路径。
CPU 对拍：\texttt{bash run.sh -r cpu -v Ascend910B4}；\textbf{1 / 2 / 8 个 Vector AI Core 三档}（见 §\ref{sec:aicore}）。
\end{abstract}

\tableofcontents
\newpage

\section{数学语义}

\subsection{输入输出}
\begin{itemize}
  \item 输入 \texttt{mat\_c\_gm}：$[16,512]$ \texttt{int32}。行 $0..7$ 为 HI 乘积行，$8..15$ 为 LO 乘积行（与 F203 Stage2 一致）。
  \item 输出 \texttt{out\_gm}：$[8,256]$ \texttt{int32}，系数 $\in [0,q)$，$q=3329$。
\end{itemize}

\subsection{RouteA（对齐 \texttt{f203\_stage3\_route\_a}）}
对每条 poly $p\in[0,7]$、输出列 $j\in[0,255]$，令偶/奇列 $c_e=2j$，$c_o=2j+1$：
\begin{align}
  hh &= C[p, c_e], \quad hl = C[p+8, c_e] \\
  lh &= C[p, c_o], \quad ll = C[p+8, c_o] \\
  raw &= hh\cdot 4096 + (hl+lh)\cdot 64 + ll \quad (\text{int64})
\end{align}

\subsection{Stage3.1（$\bmod\,q$，\texttt{stage31\_mod}）}

\textbf{本工程采用的主公式}（CANN 9.0.0，假定整除/取整语义正确）：
\begin{align}
  t &= \begin{cases}
    \lfloor raw/q \rfloor & raw \ge 0 \\
    -\lfloor (-raw)/q \rfloor & raw < 0
  \end{cases} \qquad rem = raw - q\cdot t \in [0,q)
\end{align}
对 int64 的 floor 除法，$rem$ 已落在 $[0,q)$，\textbf{数学上无需}再做 $rem\ge q$ / $rem<0$ 的归一校正。

\textbf{历史背景（ntt\_study / 交付 ONNX 写法）}：早期实验环境与部分 \texttt{Div} 向量实现存在底层正确性问题，交付图与 \texttt{mlkem\_ref.stage31\_mod} 在 $rem$ 计算后仍保留
\begin{align}
  rem &\leftarrow rem - q\cdot \mathbf{1}_{rem\ge q}, \quad
  rem \leftarrow rem + q\cdot \mathbf{1}_{rem<0}
\end{align}
作为\textbf{双校正}兜底。本仓库基于 \textbf{AscendC 9.0.0} 工具链开发，\textbf{先默认} \texttt{Div}/标量整除无已知底层错误；AscendC 标量路径仅用单次 $rem=raw-q\cdot t$。
Golden 仍调用 \texttt{stage31\_mod}（与交付 ONNX 字符串一致）；在正确整除前提下，额外两步对合法输入为\textbf{恒等}，与内核结果一致。

\section{向量实现规划}

\subsection{AI Core 规模与验证档位}
\label{sec:aicore}

\begin{longtable}{@{}L{2.8cm}L{10.5cm}@{}}
\toprule
项 & 本用例约定 \\
\midrule
核类型 & \texttt{KERNEL\_TYPE\_AIV\_ONLY}：仅使用 \textbf{Vector AI Core}（AIV），\textbf{不使用} Cube AI Core（AIC）。 \\
计划 AI Core 数 & \textbf{三档}：\texttt{aiv=1/2/8} $\Rightarrow$ \texttt{blockDim=1/2/8}（\texttt{launch\_profile.h} + \texttt{run.sh --aiv}）。 \\
命名约定 & \texttt{LAUNCH\_PROFILE=aiv=N}；\texttt{run.sh --aiv 1|2|8}（与 Stage1 一致）。 \\
数据划分（aiv=8） & 每核 1 poly：\texttt{blockIdx}$=p$；读 \texttt{mat\_c} 行 $p$、$8+p$；写 \texttt{out} 行 $p$。 \\
数据划分（aiv=2） & 每核 4 poly：\texttt{blockIdx=0,1}，核内 $p=4\cdot blockIdx..+3$ 串行；每 poly 新建 op。 \\
数据划分（aiv=1） & 单核串行 8 poly；每 poly 新建 op。 \\
每核工作量 & 每条 poly 256 系数；\texttt{kTileLen=16}、双缓冲流水（见下节）。 \\
融合算子 AIV 约束 & 含 Cube 融合时 AIV 数须为 \textbf{2 的倍数}；本用例增测 \texttt{aiv=2}（每核 4 poly）对齐最小偶数多核场景。 \\
多档验证 & \textbf{默认三档均测}：\texttt{aiv=1}、\texttt{aiv=2}、\texttt{aiv=8}；亦可 \texttt{--aiv} 单档。 \\
验证命令 & \texttt{gen\_data.py} $\rightarrow$ 各档 \texttt{ascendc\_kernels\_bbit} $\rightarrow$ \texttt{md5sum}/\texttt{cmp}/\texttt{verify\_result.py}。 \\
\bottomrule
\end{longtable}

\textbf{选型理由}：
\begin{itemize}
  \item \textbf{aiv=8}：满核 poly 并行；\textbf{aiv=2}：偶数 AIV、每核 4 poly，对齐融合 MIX 约束；\textbf{aiv=1}：单核基线。
  \item 算术含标量 int64（RouteA + mod），多档主要验证 \texttt{DataCopy} 分核与 GM 偏移正确性。
  \item 核入口与 Stage1 同构：\texttt{blockNum==1} 串行 8 poly；否则按 \texttt{8/blockNum} 均分 poly。
\end{itemize}

\subsection{分块与 Gather 策略}
\begin{itemize}
  \item \texttt{kTileLen=16}、\texttt{kPairCols=32}：每 tile 覆盖输出列 $j_0..j_0+31$（每列对应一对偶/奇 mat\_c 列）。
  \item 从 GM 连续拷贝 \texttt{int32} 对列块：HI 行偏移 $p\cdot 512 + 2j_0$，长度 $2\times\texttt{kTileLen}$（含 $hh,lh$ 交错）；
    LO 行同理得 $hl,ll$ 交错。
  \item \textbf{解交错}：在 UB 小缓冲上按索引 $2i,2i+1$ 拆出四路 \texttt{int32}——\textbf{标量}（32 次/ tile）。
  \item 备选（未在本版实现）：\texttt{DataCopyPad} + \texttt{DataCopyExtParams} 跨步搬运（07\_0265），或 \texttt{Gather} 类 ISASI API。
\end{itemize}

\section{全量 API 清单}
\label{sec:api-table}

说明：长表须用 \texttt{L\{\}} 列 + \texttt{\textbackslash apiname} 断行；\textbf{禁止}多 API 挤在同一单元格（见 Stage1 customspec）。

\begin{longtable}{@{}L{2.6cm}L{1.5cm}L{2.4cm}L{0.8cm}L{3.6cm}L{3.5cm}@{}}
\toprule
API & 分类 & 在线文档 ID & A2 & 输入/输出 dtype（本用例） & 备注 \\
\midrule
\endhead
\apiname{GlobalTensor} & 2.2 数据结构 & 07\_0007 & $\checkmark$ & GM 张量 & \texttt{SetGlobalBuffer} \\
\apiname{LocalTensor} & 2.2 数据结构 & 07\_0006 & $\checkmark$ & UB 张量 & \texttt{AllocTensor}/\texttt{GetValue} \\
\apiname{GetBlockIdx} & 2.3 系统变量 & 07\_0185 & $\checkmark$ & 核索引 & poly 划分 \\
\apiname{KERNEL\_TASK\_TYPE\_DEFAULT} & Utils / 核类型 & 见编程指南 & $\checkmark$ & \texttt{AIV\_ONLY} & 纯向量核 \\
\apiname{TPipe} & 2.3 资源管理 & 07\_0108 & $\checkmark$ & 管道 & 每 poly 新建 \\
\apiname{InitBuffer} & 2.3 资源管理 & 07\_0110 & $\checkmark$ & UB 分配 & 双缓冲队列 \\
\apiname{TQue} & 2.3 Pipe/Que & 07\_0137 & $\checkmark$ & \texttt{VECIN}/\texttt{VECOUT} & HI/LO 输入、out 输出 \\
\apiname{DataCopy} & 2.3.1 搬运 & 07\_0101 & $\checkmark$ & \texttt{int32} GM$\leftrightarrow$UB & 连续对列块 \\
\apiname{DataCopyPad} & 2.3.1 ISASI & 07\_0265 & $\checkmark$ & \texttt{int32} & \textbf{备选}跨步 gather \\
\apiname{Add} & 2.3.3 矢量 & 07\_0035 & $\checkmark$ & \texttt{int32} & $hl+lh$；见覆盖率 \\
\apiname{Muls} & 2.3.3 矢量 & 07\_0055 & $\checkmark$ & \texttt{int32} & $\times 4096,\times 64$；见覆盖率 \\
\apiname{Mul} & 2.3.3 矢量 & 07\_0037 & $\checkmark$ & \texttt{int32} & \texttt{Muls} 等价备选 \\
\apiname{Cast} & 2.3.3 类型转换 & 07\_0073 & $\checkmark$ & \texttt{int32}$\rightarrow$\texttt{int64} & 后续 int64 算术缺 \\
\apiname{Div} & 2.3.3 矢量 & 07\_0038 & $\checkmark$ & \texttt{half}/\texttt{float} & \textbf{无} \texttt{int32/int64} \\
\apiname{Compare} & 2.3.3 矢量 & 07\_0066 & $\checkmark$ & \texttt{float/half}；\texttt{int32} 仅 \texttt{EQ} & 双校正比较 \\
\apiname{CompareScalar} & 2.3.3 矢量 & 07\_0068 & $\checkmark$ & 同上 & 与 $q$、$0$ 比较 \\
\apiname{Select} & 高阶 API & 07\_0859 & $\checkmark$ & 掩码选路 & 需 \texttt{float} 链 \\
\apiname{Sub} & 2.3.3 矢量 & 07\_0036 & $\checkmark$ & \texttt{int32} & float 链 $rem-q\cdot t$ \\
\bottomrule
\end{longtable}

\section{逐步覆盖率评估}

\begin{longtable}{@{}L{0.9cm}L{3.6cm}L{3.4cm}L{1.8cm}L{3.5cm}@{}}
\toprule
步 & 数学 & 首选 API & 覆盖 & 标量/备选 \\
\midrule
G1 & 读 $C$ 对列块 & \apiname{DataCopy} & 100\% & --- \\
G2 & 拆 $hh,hl,lh,ll$ & \apiname{DataCopyPad} 跨步 & 0\%（本版） & \textbf{标量解交错} \\
R1 & $hl+lh$ & \apiname{Add} & 条件 & int32 可向量；见 R3 \\
R2 & $\times 4096,\times 64$ & \apiname{Muls} & 条件 & 乘积可超 int32 \\
R3 & $raw$ int64 合并 & \apiname{Cast}+矢量算术 & \textbf{0\%} & \textbf{标量 int64} \\
M1 & $\lfloor raw/q\rfloor$ 等 & \apiname{Div} & \textbf{0\%} & \textbf{标量}；仅 float/half \\
M2 & $rem$ 归一（历史） & \apiname{Compare}+\apiname{Select} & \textbf{0\%} & 本工程标量单次 $rem$ \\
O1 & 写 \texttt{out} & \apiname{DataCopy} & 100\% & --- \\
\bottomrule
\end{longtable}

\subsection{为何 int32 向量不能安全承担 $raw$}
单路 mat\_c 元素可达 $\sim 2\times 10^6$ 量级（256 路 int8 点积）。则 $hh\cdot 4096$ 可超过 $2^{31}$，\textbf{int32 向量乘加会溢出}；Python 使用 int64。dav\_c220 上 \texttt{Muls}/\texttt{Add} 不支持 \texttt{int64\_t}。

\subsection{Stage3.1 向量备选（未采用）}
\begin{enumerate}
  \item \textbf{float 链}：\texttt{Cast(int32$\rightarrow$float)} $\rightarrow$ \texttt{Div} $\rightarrow$ \texttt{Cast} $\rightarrow$ \texttt{Compare}/\texttt{Select}。风险：大整数除法/取整与 ONNX Div 双校正逐元素一致需范围证明。
  \item \textbf{定点乘法}：类似 \texttt{mlkem\_reduce\_to\_zq} 的移位近似（需单独验证与 \texttt{stage31\_mod} 等价）。
\end{enumerate}

\subsection{本版实现闭环结论}
\begin{itemize}
  \item \textbf{向量部分}：\texttt{DataCopy} 流水、核间划分、输出写回（$\approx$ 搬运阶段 100\%）。
  \item \textbf{标量部分}：每系数 \texttt{int64} RouteA 合并 + \texttt{Stage31ModI64}（与融合内核 \texttt{sepolyvec8\_ntt\_custom.cpp} 同语义）。
  \item \textbf{总体向量覆盖率}：算术主路径约 \textbf{30\%--40\%}（按指令条数粗估）；\textbf{功能正确性 100\%} 依赖标量缺口。
\end{itemize}

\section{实现与验证}

\begin{itemize}
  \item 内核：\texttt{f203\_stage3\_routea\_mod\_custom.cpp}（\texttt{AIV\_ONLY}，\texttt{blockDim=8}，见 §\ref{sec:aicore}）
  \item Host：\texttt{main.cpp}（\texttt{kBlockDim=8}；\texttt{ICPU\_RUN\_KF}/\texttt{f203\_stage3\_routea\_mod\_do}）
  \item Golden：\texttt{f203\_stage3\_route\_a} + \texttt{stage31\_mod}（\texttt{scripts/gen\_data.py}）
  \item 验证：\texttt{bash run.sh -r cpu -v Ascend910B4}（默认 \texttt{aiv=1/2/8} 三档）
\end{itemize}

\end{document}
